\documentclass[../main.tex]{subfiles}

\begin{document}

\ifSubfilesClassLoaded{

    %\tableofcontents
    
    \setcounter{chapter}{9}
}{}

\chapter{ HW10 }
代靖涵 25120222201319

\begin{exercise}{}{}
设二维离散随机变量$(X,Y)$的可能取值为
\[
(0,0), \quad (-1,1), \quad (-1,2), \quad (1,0),
\]
且取这些值的概率依次为$1/6, 1/3, 1/12, 5/12$, 试求$X$与$Y$各自的边际分布列.
\end{exercise}
\begin{solution}
    

\[
P\left( X= t, Y< \infty \right)= \begin{cases} \frac{5 }{12 },&t= 1 
    \\\frac{1 }{6 },&t= 0 \\ 
 \frac{1 }{3 }+ \frac{1 }{12 }= \frac{5 }{12 },&t= -1  \\ 
  0,& \text{其它}    \end{cases}  
\]

\[
P\left( X< \infty,Y= t \right)= \begin{cases} \frac{1 }{12 },&t= 2\\ 
 \frac{1 }{3 },&t= 1\\ 
  \frac{1 }{6 }+ \frac{5 }{12 }= \frac{7 }{12 },&t= 0\\ 
   0,& \text{其它}      \end{cases}  
\]

\end{solution}

\begin{exercise}{}{}
设二维随机变量$(X,Y)$的联合分布函数为
\[
F(x,y) =
\begin{cases}
1 - e^{-\lambda_1 x} - e^{-\lambda_2 y} + e^{-\lambda_1 x - \lambda_2 y - \lambda_{12} \max\{x,y\}}, & x > 0, \quad y > 0, \\
0, & \text{其他}.
\end{cases}
\]
试求$X$和$Y$各自的边际分布函数.
\end{exercise}
\begin{solution}
    
\[
F_{X}\left( x \right)= F\left( x,\infty \right)= \begin{cases} 1-e^{- \lambda _1 x},&x> 0\\ 
 0,&\text{其它} \end{cases}   
\]\begin{center}
\begin{tabular}{|c|c|c|c|}
\hline
\(X\) & \(-1\) & \(0\) & \(1\) \\
\hline
\(P\) & \(\frac{5}{12}\) & \(\frac{1}{6}\) & \(\frac{5}{12}\) \\
\hline
\end{tabular}
\end{center}
 \[
F_{Y}\left( x \right)= F\left( \infty,y \right)= \begin{cases} 1-e^{- \lambda _2 y},&y> 0\\ 
 0,&\text{其它} \end{cases}   
\]
\begin{center}
\begin{tabular}{|c|c|c|c|}
\hline
\(Y\) & \(0\) & \(1\) & \(2\) \\
\hline
\(P\) & \(\frac{7}{12}\) & \(\frac{1}{3}\) & \(\frac{1}{12}\) \\
\hline
\end{tabular}
\end{center}

\end{solution}

\begin{exercise}{}{}
求以下给出的 $(X,Y)$ 的联合密度函数的边际密度函数 $p_X(x)$ 和 $p_Y(y)$:
\begin{enumerate}
\item \[p(x,y) = \begin{cases} e^{-y}, & 0 < x < y, \\ 0, & \text{其他}; \end{cases}\]
\item \[p(x,y) = \begin{cases} \frac{5}{4}(x^2 + y), & 0 < y < 1 - x^2, \\ 0, & \text{其他}; \end{cases}\]
\item \[p(x,y) = \begin{cases} \frac{1}{x}, & 0<y<x<1, \\ 0, & \text{其他}. \end{cases}\]
\end{enumerate}
\end{exercise}

\begin{solution}
    \[
        F\left( x,y \right)= \int_{-\infty}^{x}\int_{-\infty}^{y}p\left( s,t \right)\,\mathrm{d} t\,\mathrm{d} s= \int_{-\infty}^{y}\int_{-\infty}^{x}p\left( s,t \right)\,\mathrm{d} s\,\mathrm{d} t 
        \] \[
      \begin{aligned}
        p_{X}\left( x \right)&= \frac{\mathrm{d}}{\mathrm{d}x}\lim_{y\to \infty}\int_{-\infty}^{x}\int_{-\infty}^{y}p\left( s,t \right)\,\mathrm{d} t\,\mathrm{d} s\\ 
         &=  \frac{\mathrm{d}}{\mathrm{d}x}\int_{-\infty}^{x} \int_{-\infty}^{\infty}p\left( s,t \right)\,\mathrm{d} t\,\mathrm{d} s\\ 
          &=  \int_{\mathbb{R} }p\left( x,y \right)\,\mathrm{d} y
      \end{aligned}  
        \]类似地,  \[
        p_{Y}\left( y \right)= \int_{\mathbb{R} }p\left( x,y \right)\,\mathrm{d} x  
        \]
    \begin{enumerate}
        \item  \[
        p_{X}\left( x \right)= \int_{\mathbb{R} } \chi _{\left\{ y> x \right\}}\chi _{\left\{ x> 0 \right\}}e^{-y}\,\mathrm{d} y= \begin{cases} \int_{x}^{\infty}e^{-y}\,\mathrm{d} y= e^{-x},& x> 0\\ 
         \int_{\mathbb{R} }0\,\mathrm{d} y= 0,& x\le 0 \end{cases} 
        \] \[
        p_{X}\left( x \right)= \begin{cases} e^{-x},&x> 0\\ 
         0,&x\le 0 \end{cases}  
        \] \[
        p_{Y}\left( y \right)= \int_{\mathbb{R} }\chi _{\left\{ x< y \right\}}\chi _{\left\{ x> 0 \right\}}e^{-y}\,\mathrm{d} x= \begin{cases} \int_{0}^{y}e^{-y}\,\mathrm{d} x= ye^{-y},& y> 0\\ 
         0,& y\le 0 \end{cases} 
        \] \[
        p_{Y}\left( y \right)= \begin{cases} ye^{-y},& y> 0\\ 
         0,& y\le 0 \end{cases}  
        \]
         \item  \[
        \begin{aligned}
         p_{X}\left( x \right)&= \int_{\mathbb{R} }\chi _{\left\{ y> 0 \right\}}\chi _{\left\{ 1-x^{2}> y \right\}} \frac{5 }{4 }\left( x^{2}+ y \right)\,\mathrm{d} y \\ 
          &=  \chi _{\left\{ 1-x^{2}> 0 \right\}}\int_{0}^{1-x^{2}}\frac{5 }{4 }\left( x^{2}+ y \right)\,\mathrm{d} y\\ 
           &= \chi _{\left\{ 1-x^{2} > 0\right\}} \frac{5 }{4 }\left( \left( x^{2} \right)\left( 1-x^{2} \right)+ \frac{1 }{2 }\left( 1-x^{2} \right)^{2}     \right)    \\ 
            &= \chi _{\left\{ 1-x^{2}> 0 \right\}}\frac{5 }{4 }\left( 1-x^{2} \right)\left( \frac{1 }{2 }+ \frac{1 }{2 }x^{2}   \right)   = \chi _{\left\{ 1-x^{2}> 0 \right\}}\frac{5 }{8 }\left( 1-x^{4} \right)  
        \end{aligned}
         \]因此 \[
         p_{X}\left( x \right)= \begin{cases} \frac{5 }{8 }\left( 1-x^{4} \right),& -1< x< 1\\ 
          0,& \text{其它}   \end{cases}  
         \] 
         \begin{equation}
            \begin{aligned}
                p_{Y}\left( y \right)&= \int_{\mathbb{R} }\chi _{\left\{1>  y> 0 \right\}}\chi _{\left\{ x^{2}< 1- y \right\}}\frac{5 }{4 }\left( x^{2}+ y \right)\,\mathrm{d} x\\ 
                 &= \chi _{\left\{1>  y> 0 \right\}}\int_{-\sqrt{1- y}}^{\sqrt{1- y}}   \frac{5 }{4 }\left( x^{2}+ y \right)\,\mathrm{d} x\\ 
                  &= \chi _{\left\{ 1>  y> 0 \right\}}\frac{5 }{4 }\left( \frac{2 }{3 }\left( 1- y \right)^{\frac{3}{2}}+ 2y\sqrt{1- y}   \right)  \\ 
                   &= \chi _{\left\{ 1> y> 0 \right\}}\frac{5 }{6 } \sqrt{1-y}\left(1+ 2y  \right) 
            \end{aligned}
         \end{equation}
         故 \[
         p_{Y}\left( y \right)= \begin{cases}\frac{5 }{6 } \sqrt{1-y}\left(1+ 2y  \right) ,&0< y< 1\\ 
          0,&\text{其它}    \end{cases}  
         \]
         \item  \begin{equation}
            \begin{aligned}
                p_{X}\left( x \right)&= \int_{\mathbb{R} }\frac{1 }{x }\chi _{\left\{ 0< y< x< 1 \right\}}\,\mathrm{d} y\\ 
                 &= \chi _{\left\{ 0< x< 1 \right\}}\int_{0}^{x}\frac{1 }{x }\,\mathrm{d} y\\ 
                  &= \chi _{\left\{ 0< x< 1 \right\}} 
            \end{aligned}
         \end{equation}因此 \[
         p_{X}\left( x \right)= \begin{cases} 1,&0< x< 1\\ 
          0,&\text{其它} \end{cases}  
         \]
         \begin{equation}
            \begin{aligned}
                p_{Y}\left(y \right)&= \int_{\mathbb{R} }\frac{1 }{x }\chi _{\left\{ 0< y< x< 1 \right\}}\,\mathrm{d} x\\ 
                 &=   \chi _{\left\{ 0< y< 1 \right\}}\int_{y}^{1}\frac{1 }{x }\,\mathrm{d} x\\ 
                  &= \chi _{\left\{ 0< y< 1 \right\}} \ln \frac{1 }{y } 
            \end{aligned}
         \end{equation}因此\begin{equation}
            \begin{aligned}
                p_{Y}\left( y \right)= \begin{cases} \ln \frac{1 }{y },&0< y< 1\\ 
                 0,& \text{其它}  \end{cases}  
            \end{aligned}
         \end{equation}
    \end{enumerate}
    
\end{solution}

\begin{exercise}{}{}
试验验证: 以下给出的两个不同的联合密度函数, 它们有相同的边际密度函数.
\[
p(x,y) = \begin{cases} x+y, & 0 \le x \le 1, \quad 0 \le y \le 1, \\ 0, & \text{其他}; \end{cases}
\]
\[
g(x,y) = \begin{cases} (0.5+x)(0.5+y), & 0 \le x \le 1, \quad 0 \le y \le 1, \\ 0, & \text{其他}. \end{cases}
\]
\end{exercise}

\begin{proof}
    \begin{equation}
        \begin{aligned}
            p_{X}\left( x \right)&=  \int_{\mathbb{R} }\chi _{\left\{ 0\le x\le 1 \right\}}\chi _{\left\{ 0\le y\le 1 \right\}}\left( x+ y \right)\,\mathrm{d} y  \\ 
             &= \chi _{\left\{ 0\le x\le 1 \right\}}\int_{0}^{1}\left( x+ y \right)\,\mathrm{d} y\\ 
              &= \chi _{\left\{ 0\le x\le 1 \right\}}\left( x+ \frac{1 }{2 }  \right)  
        \end{aligned}
    \end{equation}
    \begin{equation}
        \begin{aligned}
            p_{Y}\left( y \right)&= \int_{\mathbb{R} }\chi _{\left\{ 0\le y\le 1 \right\}} \chi _{0\le x\le 1}\left( x+ y \right)\,\mathrm{d} x\\ 
             &=  \chi _{\left\{ 0\le y\le 1 \right\}}\int_{0}^{1}\left( x+ y \right)\,\mathrm{d} x \\ 
              &= \chi _{\left\{ 0\le y\le 1 \right\}}\left( y+ \frac{1 }{2 }  \right) 
        \end{aligned}
    \end{equation}
\begin{equation}
    \begin{aligned}
          g_{X}\left( x \right)&= \int_{\mathbb{R} }\chi _{\left\{ 0\le x\le 1 \right\}}\chi _{\left\{ 0\le y\le 1 \right\}}\left(\frac{1 }{2 }+ x  \right)  \left( \frac{1 }{2 }+ y  \right)  \,\mathrm{d} y\\ 
           &= \chi _{\left\{ 0\le x\le 1 \right\}}\left( x+ \frac{1 }{2 }  \right) \int_{0}^{1}\left( y+ \frac{1 }{2 }  \right)\,\mathrm{d} y \\ 
            &= \chi _{\left\{ 0\le x\le 1 \right\}}\left( x+ \frac{1 }{2 }  \right) 
    \end{aligned}
\end{equation}

\begin{equation}
    \begin{aligned}
        g_{Y}\left( y \right)&= \int_{\mathbb{R} }\chi _{\left\{ 0\le x\le 1 \right\}}\chi _{\left\{ 0\le y\le 1 \right\}}\left( \frac{1 }{2 }+ x  \right)\left( \frac{1 }{2 }+ y  \right)\,\mathrm{d} x\\ 
         &= \chi _{\left\{ 0\le y\le 1 \right\}}\left( y+ \frac{1 }{2 }  \right)\int_{0}^{1}\left( x+ \frac{1 }{2 }  \right)\,\mathrm{d} x\\ 
          &= \chi _{\left\{ 0\le y\le 1 \right\}}\left( y+ \frac{1 }{2 }  \right)      
    \end{aligned}
\end{equation}
\end{proof}因此 \[
p_{X}= g_{X},\quad p_{Y}= g_{Y}
\]
\end{document}